\subsection{补算：流函数、磁通函数与Poisson括号的三个符号}
本节沿用$\B_\perp=\hat z\times\grad\psi$、$\vel_\perp=\hat z\times\grad\phi$。直接叉乘：
\begin{equation}
 B_x=-\partial_y\psi,\quad B_y=\partial_x\psi,\qquad
 u_x=-\partial_y\phi,\quad u_y=\partial_x\phi.
\end{equation}
因此
\begin{align}
 \vel_\perp\cdot\grad g
 &=-\partial_y\phi\,\partial_xg+\partial_x\phi\,\partial_yg
 =[\phi,g],\\
 (\grad\times\vel)_z
 &=\partial_xu_y-\partial_yu_x=\nabla_\perp^2\phi,\\
 (\grad\times\B)_z&=\nabla_\perp^2\psi.
\end{align}
又因通常$\B_\perp=\grad A_z\times\hat z$，这里$A_z=-\psi$。取二维势规范后$E_z=-\partial_tA_z=\partial_t\psi$。如果使用相反的磁通定义，三个符号应一起更换；只改其中一个会改变能量守恒和不稳定性。

\subsection{极化漂移的符号：从离子垂直动量再算一次}
设$\B=B_g\hat z$，最低阶漂移$\vel_E=\E_\perp\times\B/B_g^2$，满足
$\E_\perp+\vel_E\times\B=0$。
将$\vel_{i\perp}=\vel_E+\vel_{\rm pol}$代入冷离子垂直动量，在下一阶保留
\begin{equation}
 m_iD_E\vel_E=e\vel_{\rm pol}\times B_g\hat z.
\end{equation}
两边左乘$\hat z\times$，用$\hat z\times(\vel_{\rm pol}\times\hat z)=\vel_{\rm pol}$，得
\begin{equation}
 \vel_{\rm pol}={1\over\Omega_{ci}}\hat z\times D_E\vel_E.
\end{equation}
在均匀导引场、静电势$V$且保留相应约化阶次时，
$\vel_E=\hat z\times\grad V/B_g$，于是
$\vel_{\rm pol}=-D_E\grad V/(B_g\Omega_{ci})$。
把其散度放入连续性，才出现密度与涡量的关系。一般$D_E$与梯度并不交换，但在此均匀导引场、二维静电$E\times B$流的特定条件下，
\begin{equation}
 \grad\cdot(D_E\grad V)
 =D_E\nabla^2V+(\partial_i u_{E,j})(\partial_j\partial_iV).
\end{equation}
最后一项按$x,y$展开是
$(-V_{xy}V_{xx}+V_{xx}V_{xy}-V_{yy}V_{xy}+V_{xy}V_{yy})/B_g=0$。
在小密度扰动的所保留阶次，连续性因此给
\begin{equation}
 D_EN=-\grad\cdot\vel_{\rm pol}
 ={1\over B_g\Omega_{ci}}D_E\nabla^2V
 =\rho_{\rm sound}^2D_E\nabla^2\Phi.
\end{equation}
这就恢复了被输运的密度--涡量组合。若导引场不均匀、允许更强压缩或扩展到更高阶，就要重算这些抵消。电子平行动量的抗磁漂移惯性还需要与同阶压力修正一起处理，不能把这里的离子漂移计算当作完整非线性双流体到两场模型的全部证明。\citep{Schep1994,L049}

\subsection{初始密度模为什么决定能否只解两个场}
记$C_N=N-\rho_{\rm sound}^2\nabla_\perp^2\Phi$。约化离子连续性给
$D_EC_N=0$。若$t=0$时$C_N=0$且边界不注入该量，则沿每条$E\times B$流线一直为零，因此可以用
$N=\rho_{\rm sound}^2\nabla_\perp^2\Phi$代替密度方程。

若初始$C_N$是非零函数，它只会被输运，不会因“想使用两场模型”而消失。此时消去密度后会多出与$C_N$有关的项，或者应独立推进第三个场。这是模型约化中很典型的一种条件：除了小参数，还需要允许的初始数据与边界数据。

\subsection{两个广义磁通的不变量：展示所有抵消项}
为缩短式子，以下写$\rho_*=\rho_{\rm sound}$、$j=\nabla_\perp^2\psi$，
$D_\phi=\partial_t+[\phi,\cdot]$。两场方程为
\begin{equation}
 D_\phi F=\rho_*^2[\omega,\psi],\qquad
 D_\phi\omega=[\psi,j],\qquad F=\psi-d_e^2j.
\end{equation}
先试构造$G=F+a\omega$，希望存在流函数$\phi+b\psi$使其被纯输运。计算
\begin{align}
 \partial_tG+[\phi+b\psi,G]
 &=\rho_*^2[\omega,\psi]+a[\psi,j]+b[\psi,F]+ab[\psi,\omega]\\
 &=(\rho_*^2-ab)[\omega,\psi]+(a-bd_e^2)[\psi,j],
\end{align}
其中用了$[\psi,F]=-d_e^2[\psi,j]$。要让任意场的两项都消失，必须有
$ab=\rho_*^2$、$a=bd_e^2$，所以
\begin{equation}
 a=\pm d_e\rho_*,\quad b=\pm{\rho_*\over d_e},\qquad
 G_\pm=F\pm d_e\rho_*\omega.
\end{equation}
这不是凭空猜出两个组合，而是解了两个系数条件。$d_e=0$或$\rho_*=0$的退化极限要回到原方程处理，不能对含$1/d_e$的表达式直接代零。

\subsection{能量守恒的逐项验算}
在周期或足以消去边界通量的条件下，分部积分给
$\int f[g,h]=\int g[h,f]=\int h[f,g]$，并有$\int f[g,f]=0$。
两场能量的时间导数可以先整理成
\begin{equation}
 {\dd\mathcal H\over\dd t}
 =\int\{-j\,\partial_tF+(-\phi+\rho_*^2\omega)\partial_t\omega\}\dd^2x.
\end{equation}
例如磁能部分
$\int\grad\psi\cdot\grad\partial_t\psi=-\int j\partial_t\psi$；
电子惯性部分是$d_e^2\int j\,\partial_tj$，两者合成$-\int j\partial_tF$。
流动能部分用$\omega=\nabla^2\phi$得到$-\int\phi\partial_t\omega$。

代入两场方程，六类项成为
\begin{align}
 {\dd\mathcal H\over\dd t}
 =\int\{&
 j[\phi,F]-\rho_*^2j[\omega,\psi]
 +\phi[\phi,\omega]-\rho_*^2\omega[\phi,\omega]\\
 &-\phi[\psi,j]+\rho_*^2\omega[\psi,j]\}\dd^2x.\nonumber
\end{align}
第三、四项各自为零。第一项中的$-d_e^2j[\phi,j]$也为零，余下的$\int j[\phi,\psi]$与第五项抵消；第二项与第六项由循环恒等式抵消。于是$\dd\mathcal H/\dd t=0$。

验算保留了磁能、电子漂移能、垂直流动能和密度/极化相关能。若人为删除其中某项，或给Poisson括号写错一个符号，这些抵消通常就不成立。这是比“能量曲线看起来平滑”更直接的模型检查。
